\documentclass[10pt,a4paper]{amsart}
\usepackage{amsmath,amssymb,amsthm,graphicx}
% For commutative diagrams we use Xy-pic
\usepackage[all]{xy}

% We use 2cell for 2-commutative diagrams.
\xyoption{2cell}
% Theorem environments.
%
\theoremstyle{plain}
\newtheorem{theorem}[subsection]{Theorem}
\newtheorem{proposition}[subsection]{Proposition}
\newtheorem{lemma}[subsection]{Lemma}

\theoremstyle{definition}
\newtheorem{definition}[subsection]{Definition}
\newtheorem{example}[subsection]{Example}
\newtheorem{exercise}[subsection]{Exercise}
\newtheorem{situation}[subsection]{Situation}

\theoremstyle{remark}
\newtheorem{remark}[subsection]{Remark}
\newtheorem{remarks}[subsection]{Remarks}

\numberwithin{equation}{subsection}

% Macros
%
\def\lim{\mathop{\mathrm{lim}}\nolimits}
\def\colim{\mathop{\mathrm{colim}}\nolimits}
\def\Spec{\mathop{\mathrm{Spec}}}
\def\Hom{\mathop{\mathrm{Hom}}\nolimits}
\def\Ext{\mathop{\mathrm{Ext}}\nolimits}
\def\SheafHom{\mathop{\mathcal{H}\!\mathit{om}}\nolimits}
\def\SheafExt{\mathop{\mathcal{E}\!\mathit{xt}}\nolimits}
\def\Sch{\mathit{Sch}}
\def\Mor{\mathop{\mathrm{Mor}}\nolimits}
\def\Ob{\mathop{\mathrm{Ob}}\nolimits}
\def\Sh{\mathop{\mathit{Sh}}\nolimits}
\def\NL{\mathop{N\!L}\nolimits}
\def\CH{\mathop{\mathrm{CH}}\nolimits}
\def\proetale{{pro\text{-}\acute{e}tale}}
\def\etale{{\acute{e}tale}}
\def\QCoh{\mathit{QCoh}}
\def\Ker{\mathop{\mathrm{Ker}}}
\def\Im{\mathop{\mathrm{Im}}}
\def\Coker{\mathop{\mathrm{Coker}}}
\def\Coim{\mathop{\mathrm{Coim}}}

% Boxtimes
%
\DeclareMathSymbol{\boxtimes}{\mathbin}{AMSa}{"02}

%
% Macros for moduli stacks/spaces
%
\def\QCohstack{\mathcal{QC}\!\mathit{oh}}
\def\Cohstack{\mathcal{C}\!\mathit{oh}}
\def\Spacesstack{\mathcal{S}\!\mathit{paces}}
\def\Quotfunctor{\mathrm{Quot}}
\def\Hilbfunctor{\mathrm{Hilb}}
\def\Curvesstack{\mathcal{C}\!\mathit{urves}}
\def\Polarizedstack{\mathcal{P}\!\mathit{olarized}}
\def\Complexesstack{\mathcal{C}\!\mathit{omplexes}}
% \Pic is the operator that assigns to X its picard group, usage \Pic(X)
% \Picardstack_{X/B} denotes the Picard stack of X over B
% \Picardfunctor_{X/B} denotes the Picard functor of X over B
\def\Pic{\mathop{\mathrm{Pic}}\nolimits}
\def\Picardstack{\mathcal{P}\!\mathit{ic}}
\def\Picardfunctor{\mathrm{Pic}}
\def\Deformationcategory{\mathcal{D}\!\mathit{ef}}
\usepackage{cleveref}
\crefname{originaltheorem}{original theorem}{original theorems}
\Crefname{originaltheorem}{Original theorem}{Original theorems}
\makeatletter
\newlabel{theorem-cech-ss}{{03OW}{source}{}{sourcepointer-1}{}}
\newlabel{theorem-cech-ss@cref}{{[originaltheorem][0][]03OW}{[1][0][]source}{}{}{}}
\makeatother
\crefname{originaltheorem}{original theorem}{original theorems}
\Crefname{originaltheorem}{Original theorem}{Original theorems}
\makeatletter
\newlabel{theorem-picard-group}{{03P8}{source}{}{sourcepointer-2}{}}
\newlabel{theorem-picard-group@cref}{{[originaltheorem][0][]03P8}{[1][0][]source}{}{}{}}
\makeatother
\crefname{originallemma}{original lemma}{original lemmas}
\Crefname{originallemma}{Original lemma}{Original lemmas}
\makeatletter
\newlabel{lemma-kummer-sequence}{{03PL}{source}{}{sourcepointer-3}{}}
\newlabel{lemma-kummer-sequence@cref}{{[originallemma][0][]03PL}{[1][0][]source}{}{}{}}
\makeatother
\crefname{originallemma}{original lemma}{original lemmas}
\Crefname{originallemma}{Original lemma}{Original lemmas}
\makeatletter
\newlabel{lemma-kummer-sequence-syntomic}{{040N}{source}{}{sourcepointer-4}{}}
\newlabel{lemma-kummer-sequence-syntomic@cref}{{[originallemma][0][]040N}{[1][0][]source}{}{}{}}
\makeatother
\crefname{originalsection}{original section}{original sections}
\Crefname{originalsection}{Original section}{Original sections}
\makeatletter
\newlabel{sites-modules-section-internal-hom}{{04TT}{source}{}{sourcepointer-5}{}}
\newlabel{sites-modules-section-internal-hom@cref}{{[originalsection][0][]04TT}{[1][0][]source}{}{}{}}
\makeatother
\crefname{originallemma}{original lemma}{original lemmas}
\Crefname{originallemma}{Original lemma}{Original lemmas}
\makeatletter
\newlabel{sites-cohomology-lemma-torsors-h1}{{03AJ}{source}{}{sourcepointer-6}{}}
\newlabel{sites-cohomology-lemma-torsors-h1@cref}{{[originallemma][0][]03AJ}{[1][0][]source}{}{}{}}
\makeatother
\crefname{originallemma}{original lemma}{original lemmas}
\Crefname{originallemma}{Original lemma}{Original lemmas}
\makeatletter
\newlabel{sites-cohomology-lemma-trivial-torsor}{{03AI}{source}{}{sourcepointer-7}{}}
\newlabel{sites-cohomology-lemma-trivial-torsor@cref}{{[originallemma][0][]03AI}{[1][0][]source}{}{}{}}
\makeatother
\title{First cohomology with mu\_n coefficients as pairs of an invertible sheaf with a trivialized n-th tensor power}
\author{The Stacks Project Authors}
\date{stacks-project:etale-cohomology.tex:local-snapshot-20261011; GNU FDL 1.2 or later}
\begin{document}
\maketitle
\noindent\textit{Editorial note.} Counted proposition: Tag 040Q, etale-cohomology.tex lines 3056-3069 with its own proof at lines 3071-3178 (statement immediately followed by its proof; no non-adjacent pairing).\par\medskip
\noindent\textit{Editorial note.} The two context blocks reproduce the author's own definition of the pairs group (including the displayed description 040P of its isomorphisms) and the Kummer long exact sequence that the proof invokes as 'the long exact cohomology sequence above'. They are supporting material, not additional counted problems.\par\medskip
\noindent\textit{Editorial note.} The Cech-to-cohomology spectral sequence (Tag 03OW) is counted as a separate task of this batch; here it appears only as an external tag pointer.\par\medskip
\noindent\textit{Editorial note.} External references are navigation pointers to stable Stacks tags; the two constructions of the class and the surjectivity argument are retained in full.\par\medskip
\noindent\textit{Allowed standard prerequisites (editorial).}\begin{itemize}
\item Cech-to-cohomology spectral sequence on a site. Stacks tag 03OW (source lines 1770-1830; counted as a separate task of this batch). Maps the Cech class of the covering built in the proof to H\textasciicircum{}1\_syntomic(S, mu\_n) (author's lines 3100-3106).
\item Hilbert 90: H\textasciicircum{}1 of the multiplicative group is the Picard group. Stacks tag 03P8 (source lines 2497-2508). Turns the image of xi in H\textasciicircum{}1\_syntomic(S, G\_m) into an invertible sheaf whose n-th tensor power is trivial (author's lines 3156-3159).
\item The Kummer sequence in the syntomic and fppf topologies. Stacks tags 040N and 03PL (source lines 2921 and 2842). Gives the long exact sequence of lines 2975-2997 and the exactness that lets the author write xi - xi' = partial(f) at lines 3163-3165 and, for n invertible, replace the syntomic covering by an etale one at lines 3172-3177.
\item Torsors on a site and their cohomology classes. Stacks tags 03AJ and 03AI; Sites-Modules, Section 04TT (internal hom sheaves). Underlies the second construction of the class as the mu\_n-torsor of isomorphisms (lines 3117-3129) and the criterion that a torsor with a global section is trivial (lines 3147-3149).
\end{itemize}\medskip
\ifdefined\hypertarget\hypertarget{sourcepointer-1}{}\fi
\noindent\textit{Source reference.} Original theorem 03OW: sources/stacks-etale-cohomology/pinned-original-source.tex, source lines 1771--1771. Navigation pointer to the Cech-to-cohomology spectral sequence, cited by the author at line 3105; the index 03OW is also counted as a separate task of this batch and its proof is not reproduced here.\par\medskip
\ifdefined\hypertarget\hypertarget{sourcepointer-2}{}\fi
\noindent\textit{Source reference.} Original theorem 03P8: sources/stacks-etale-cohomology/pinned-original-source.tex, source lines 2498--2498. Hilbert 90, cited at line 3157 to turn a class in H\textasciicircum{}1(., G\_m) into an invertible sheaf; a foundational input, not a substitute for the counted surjectivity argument.\par\medskip
\ifdefined\hypertarget\hypertarget{sourcepointer-3}{}\fi
\noindent\textit{Source reference.} Original lemma 03PL: sources/stacks-etale-cohomology/pinned-original-source.tex, source lines 2842--2842. Exactness of the Kummer sequence, cited at lines 3175 and 3177 for the last step of the proof and for the etale case.\par\medskip
\ifdefined\hypertarget\hypertarget{sourcepointer-4}{}\fi
\noindent\textit{Source reference.} Original lemma 040N: sources/stacks-etale-cohomology/pinned-original-source.tex, source lines 2921--2921. Syntomic Kummer sequence, cited in the retained context at line 2979 to produce the long exact sequence displayed at lines 2982-2995.\par\medskip
\ifdefined\hypertarget\hypertarget{sourcepointer-5}{}\fi
\noindent\textit{Source reference.} Original section 04TT: sources/stacks-etale-cohomology/tags.txt, source lines 6266--6266. Pointer for the internal hom sheaf into which the presheaf mu\_n(L, alpha) is embedded at lines 3117-3120.\par\medskip
\ifdefined\hypertarget\hypertarget{sourcepointer-6}{}\fi
\noindent\textit{Source reference.} Original lemma 03AJ: sources/stacks-etale-cohomology/tags.txt, source lines 4312--4312. Standard correspondence between torsors and H\textasciicircum{}1 classes, cited at line 3128 for the second construction.\par\medskip
\ifdefined\hypertarget\hypertarget{sourcepointer-7}{}\fi
\noindent\textit{Source reference.} Original lemma 03AI: sources/stacks-etale-cohomology/tags.txt, source lines 4311--4311. Criterion that a torsor with a global section is trivial, cited at line 3149 for item (4) of the author's verification list.\par\medskip
\noindent
By
Theorem \ref{theorem-picard-group}
and
Lemma \ref{lemma-kummer-sequence-syntomic}
and general properties of cohomology we obtain
the long exact cohomology sequence
$$
\xymatrix{
0 \ar[r] &
H_{fppf}^0(S, \mu_{n, S}) \ar[r] &
\Gamma(S, \mathcal{O}_S^*) \ar[r]^{(\cdot)^n} &
\Gamma(S, \mathcal{O}_S^*) \ar@(rd, ul)[rdllllr]
\\
& H_{fppf}^1(S, \mu_{n, S}) \ar[r] &
\Pic(S) \ar[r]^{(\cdot)^n} &
\Pic(S) \ar@(rd, ul)[rdllllr] \\
& H_{fppf}^2(S, \mu_{n, S}) \ar[r] &
\ldots
}
$$
for any scheme $S$ and any integer $n$. Of course there is a similar sequence
with syntomic cohomology.

\medskip\noindent
Let $n \in \mathbf{N}$ and let $S$ be any scheme.
There is another more direct way to describe the first cohomology group with
values in $\mu_n$. Consider pairs
$(\mathcal{L}, \alpha)$ where $\mathcal{L}$ is an invertible sheaf on $S$
and $\alpha : \mathcal{L}^{\otimes n} \to \mathcal{O}_S$ is a trivialization
of the $n$th tensor power of $\mathcal{L}$.
Let $(\mathcal{L}', \alpha')$ be a second such pair.
An isomorphism $\varphi : (\mathcal{L}, \alpha) \to (\mathcal{L}', \alpha')$
is an isomorphism $\varphi : \mathcal{L} \to \mathcal{L}'$ of invertible
sheaves such that the diagram
$$
\xymatrix{
\mathcal{L}^{\otimes n} \ar[d]_{\varphi^{\otimes n}} \ar[r]_\alpha &
\mathcal{O}_S \ar[d]^1 \\
(\mathcal{L}')^{\otimes n} \ar[r]^{\alpha'} &
\mathcal{O}_S \\
}
$$
commutes. Thus we have
\begin{equation}
\label{equation-isomorphisms-pairs}
\mathit{Isom}_S((\mathcal{L}, \alpha), (\mathcal{L}', \alpha'))
=
\left\{
\begin{matrix}
\emptyset & \text{if} & \text{they are not isomorphic} \\
H^0(S, \mu_{n, S})\cdot \varphi & \text{if} &
\varphi \text{ isomorphism of pairs}
\end{matrix}
\right.
\end{equation}
Moreover, given two pairs $(\mathcal{L}, \alpha)$, $(\mathcal{L}', \alpha')$
the tensor product
$$
(\mathcal{L}, \alpha) \otimes (\mathcal{L}', \alpha')
=
(\mathcal{L} \otimes \mathcal{L}', \alpha \otimes \alpha')
$$
is another pair. The pair $(\mathcal{O}_S, 1)$ is an identity for this
tensor product operation, and an inverse is given by
$$
(\mathcal{L}, \alpha)^{-1} = (\mathcal{L}^{\otimes -1}, \alpha^{\otimes -1}).
$$
Hence the collection of isomorphism classes of pairs forms an abelian group.
Note that
$$
(\mathcal{L}, \alpha)^{\otimes n}
=
(\mathcal{L}^{\otimes n}, \alpha^{\otimes n})
\xrightarrow{\alpha}
(\mathcal{O}_S, 1)
$$
is an isomorphism
hence every element of this group has order dividing $n$. We warn the reader
that this group is in general {\bf not} the $n$-torsion in $\Pic(S)$.

\begin{lemma}
\label{lemma-describe-h1-mun}
Let $S$ be a scheme. There is a canonical identification
$$
H_\etale^1(S, \mu_n) =
\text{group of pairs }(\mathcal{L}, \alpha)\text{ up to isomorphism as above}
$$
if $n$ is invertible on $S$. In general we have
$$
H_{fppf}^1(S, \mu_n) =
\text{group of pairs }(\mathcal{L}, \alpha)\text{ up to isomorphism as above}.
$$
The same result holds with fppf replaced by syntomic.
\end{lemma}

\begin{proof}
We first prove the second isomorphism.
Let $(\mathcal{L}, \alpha)$ be a pair as above.
Choose an affine open covering $S = \bigcup U_i$ such that
$\mathcal{L}|_{U_i} \cong \mathcal{O}_{U_i}$. Say $s_i \in \mathcal{L}(U_i)$
is a generator. Then $\alpha(s_i^{\otimes n}) = f_i \in \mathcal{O}_S^*(U_i)$.
Writing $U_i = \Spec(A_i)$ we see there exists a global
relative complete intersection $A_i \to B_i = A_i[T]/(T^n - f_i)$
such that $f_i$ maps to an $n$th power in $B_i$. In other words, setting
$V_i = \Spec(B_i)$ we obtain a syntomic covering
$\mathcal{V} = \{V_i \to S\}_{i \in I}$ and trivializations
$\varphi_i : (\mathcal{L}, \alpha)|_{V_i} \to (\mathcal{O}_{V_i}, 1)$.

\medskip\noindent
We will use this result (the existence of the covering $\mathcal{V}$)
to associate to this pair a cohomology class in
$H^1_{syntomic}(S, \mu_{n, S})$. We give two (equivalent) constructions.

\medskip\noindent
First construction: using {\v C}ech cohomology.
Over the double overlaps $V_i \times_S V_j$ we have the isomorphism
$$
(\mathcal{O}_{V_i \times_S V_j}, 1)
\xrightarrow{\text{pr}_0^*\varphi_i^{-1}}
(\mathcal{L}|_{V_i \times_S V_j}, \alpha|_{V_i \times_S V_j})
\xrightarrow{\text{pr}_1^*\varphi_j}
(\mathcal{O}_{V_i \times_S V_j}, 1)
$$
of pairs. By (\ref{equation-isomorphisms-pairs}) this is given by an
element $\zeta_{ij} \in \mu_n(V_i \times_S V_j)$. We omit the verification
that these $\zeta_{ij}$'s give a $1$-cocycle, i.e., give
an element $(\zeta_{i_0i_1}) \in \check C(\mathcal{V}, \mu_n)$
with $d(\zeta_{i_0i_1}) = 0$. Thus its class is an element in
$\check H^1(\mathcal{V}, \mu_n)$ and by
Theorem \ref{theorem-cech-ss}
it maps to a cohomology class in $H^1_{syntomic}(S, \mu_{n, S})$.

\medskip\noindent
Second construction: Using torsors. Consider the presheaf
$$
\mu_n(\mathcal{L}, \alpha) :
U
\longmapsto
\mathit{Isom}_U((\mathcal{O}_U, 1), (\mathcal{L}, \alpha)|_U)
$$
on $(\Sch/S)_{syntomic}$.
We may view this as a subpresheaf of
$\SheafHom_\mathcal{O}(\mathcal{O}, \mathcal{L})$ (internal hom
sheaf, see
Modules on Sites, Section \ref{sites-modules-section-internal-hom}).
Since the conditions defining this subpresheaf are local, we see that it is
a sheaf.
By (\ref{equation-isomorphisms-pairs}) this sheaf has a free action of
the sheaf $\mu_{n, S}$. Hence the only thing we have to check is that
it locally has sections. This is true because of the existence of the
trivializing cover $\mathcal{V}$. Hence $\mu_n(\mathcal{L}, \alpha)$
is a $\mu_{n, S}$-torsor and by
Cohomology on Sites, Lemma \ref{sites-cohomology-lemma-torsors-h1}
we obtain a corresponding element of $H^1_{syntomic}(S, \mu_{n, S})$.

\medskip\noindent
Ok, now we have to still show the following
\begin{enumerate}
\item The two constructions give the same cohomology class.
\item Isomorphic pairs give rise to the same cohomology class.
\item The cohomology class of
$(\mathcal{L}, \alpha) \otimes (\mathcal{L}', \alpha')$
is the sum of the cohomology classes of
$(\mathcal{L}, \alpha)$ and $(\mathcal{L}', \alpha')$.
\item If the cohomology class is trivial, then the pair is trivial.
\item Any element of $H^1_{syntomic}(S, \mu_{n, S})$ is the
cohomology class of a pair.
\end{enumerate}
We omit the proof of (1). Part (2) is clear from the second construction,
since isomorphic torsors give the same cohomology classes.
Part (3) is clear from the first construction, since the resulting
{\v C}ech classes add up. Part (4) is clear from the second construction
since a torsor is trivial if and only if it has a global section, see
Cohomology on Sites, Lemma \ref{sites-cohomology-lemma-trivial-torsor}.

\medskip\noindent
Part (5) can be seen as follows (although a direct proof would be
preferable). Suppose $\xi \in H^1_{syntomic}(S, \mu_{n, S})$.
Then $\xi$ maps to an element
$\overline{\xi} \in H^1_{syntomic}(S, \mathbf{G}_{m, S})$
with $n \overline{\xi} = 0$. By
Theorem \ref{theorem-picard-group}
we see that $\overline{\xi}$ corresponds to an invertible sheaf $\mathcal{L}$
whose $n$th tensor power is isomorphic to $\mathcal{O}_S$.
Hence there exists a pair $(\mathcal{L}, \alpha')$ whose cohomology
class $\xi'$ has the same image $\overline{\xi'}$ in
$H^1_{syntomic}(S, \mathbf{G}_{m, S})$. Thus it suffices to show
that $\xi - \xi'$ is the class of a pair. By construction, and the
long exact cohomology sequence above, we see that
$\xi - \xi' = \partial(f)$ for some $f \in H^0(S, \mathcal{O}_S^*)$.
Consider the pair $(\mathcal{O}_S, f)$. We omit the verification
that the cohomology class of this pair is $\partial(f)$, which
finishes the proof of the first identification (with fppf replaced
with syntomic).

\medskip\noindent
To see the first, note that if $n$ is invertible on $S$, then the
covering $\mathcal{V}$ constructed in the first part of the proof
is actually an \'etale covering (compare with the proof of
Lemma \ref{lemma-kummer-sequence}). The rest of the proof is independent
of the topology, apart from the very last argument which uses that
the Kummer sequence is exact, i.e., uses Lemma \ref{lemma-kummer-sequence}.
\end{proof}

\end{document}
